\section{Evaluation}
\label{sec:evaluation}

\subsection{Experimental Setup}

All experiments are conducted on the single-VM deployment described in
Section~\ref{sec:architecture}. Each threat demo is executed in isolation after
a full reset: Wazuh OpenSearch indices are wiped, Docker volumes are cleared, and
all log files are deleted. This guarantees that counter values and rule hit counts
reflect a single clean run. The attacker container is a separate Incus instance
with its own IP, so blocking it does not affect access to Grafana or Prometheus
from the host.

Each demo script follows a three-phase structure: a warm-up period to seed
Prometheus counters, the attack phase, and a post-attack observation window.
Grafana panels are read at the end of each run to confirm detection and response.

\subsection{Threshold Analysis}

The detection thresholds were chosen for the simulated environment and require
evaluation against realistic operational assumptions.

\begin{itemize}[leftmargin=*]
  \item \textbf{DoS - 50 requests per 10 seconds:} this is achievable by a
        browser with parallel resource loading, so it is on the aggressive side
        for a production portal. In deployment, the threshold should be calibrated
        against baseline traffic measurements from the specific operator population.
  \item \textbf{APT recon - 5 unknown paths in 30 seconds:} realistic for
        deliberate enumeration scripts. A slow scanner probing one endpoint per
        minute would not trigger rule 100203 and could enumerate indefinitely
        without detection.
  \item \textbf{APT honeypot - 3 hits in 5 minutes:} permissive. A single
        accidental access to a honeypot endpoint already constitutes an anomaly
        of note. A threshold of 2 or even 1 hit would reduce the window for
        credential extraction.
  \item \textbf{Alert flood - 10 emails per 60 seconds:} the platform has 3 zones
        with 3 sensor types each (AQI, temperature, noise), yielding up to 9
        simultaneous threshold crossings in a genuine multi-zone incident. A burst
        of 9 legitimate alerts falls just below the threshold. A burst of 10
        crosses it. This boundary is narrow and would require adjustment in a
        deployment with more zones or sensor types.
  \item \textbf{Postfix rate limit - 30 messages per minute:} never triggered
        in the current demo setup and not realistic as a standalone limit. It
        serves only as a failsafe and should be sized to the actual expected alert
        volume.
  \item \textbf{SpamAssassin threshold - 8 stars:} the margin between legitimate
        emails (4-5) and drug spam (10-16) is comfortable in testing. The risk
        is at the boundary: a legitimately composed alert with unusual keyword
        combinations could approach 8, while a carefully obfuscated spam message
        could score below it.
\end{itemize}

\subsection{Scalability}

The zone aggregator design provides a structural bandwidth advantage: each
aggregator batches readings from 3 sensors and forwards a single payload every
10 seconds to the ingestion service. As the number of sensors per zone grows,
the reduction in upstream traffic scales proportionally. For $K$ sensors per zone,
the upstream message rate is $1/K$ of the per-sensor rate.

The Wazuh deployment does not scale horizontally in the current implementation.
A single agent monitors all log files, and all 13 rules are evaluated on one
host. In a real Smart City deployment with tens of zones and distributed services,
a Wazuh cluster with per-zone agents and a centralised manager would be required.
The structured JSON log format and the \texttt{event\_type} field convention
adopted throughout the platform are forward-compatible with this architecture.

Prometheus scraping at 15-second intervals (5 seconds for the ops portal and Wazuh
exporter) is adequate for the current service count. As the number of services
grows, scrape duration would increase. A federated Prometheus setup or remote
write to a time-series database would be needed to avoid staleness.

\subsection{Limitations}

\textbf{APT scope.} The implemented attack chain demonstrates reconnaissance and
honeypot probing but does not proceed beyond credential capture. A complete APT
lifecycle would include establishing a persistent backdoor or reverse shell, using
the stolen token to exfiltrate data from all three zones, and pivoting to
command-and-control infrastructure. These stages are not implemented. The demo
stops at the point where the \texttt{INTERNAL\_TOKEN} is extracted.

\textbf{Mail server trust model.} The Postfix relay currently accepts inbound mail
from any external sender and outbound mail from any internal service. In production,
the relay should be restricted to accept alert emails only from the alert service's
internal IP. Without this restriction, any compromised internal container could
send email through the relay, bypassing the circuit breaker entirely.

\textbf{Single-IP blocking.} The active response blocks the attacking IP for 300
seconds before automatically lifting the block. There is no persistent ban list; a
patient attacker can resume after each cooldown. Additionally, single-IP blocking
is ineffective against a large botnet where each bot sends traffic below the
individual threshold, as noted in Section~\ref{sec:botnet}.

\textbf{No TLS on inter-service communication.} Service-to-service calls use
plain HTTP. The \texttt{x-internal-token} header is therefore transmitted in
cleartext and could be intercepted on a compromised network segment. In a
production environment, TLS with mutual certificate authentication would replace
the shared-token scheme.

\textbf{Active response conservatism.} IP blocking was chosen over more
disruptive responses such as service termination or container isolation.
More aggressive actions would cause downtime for legitimate users sharing
the affected service and would destroy in-flight state and log entries that
are valuable for post-incident forensic analysis. IP blocking preserves service
availability for all other clients and keeps the full log record intact.

\textbf{Rule threshold validation.} All detection thresholds were tuned against
the simulated attack scripts only. No baseline traffic capture from legitimate
operator behaviour was available to validate that the chosen thresholds do not
produce false positives under realistic load. In a production deployment, a
silent monitoring period with rules in log-only mode would be required before
active response is enabled.

\textbf{Single-host lab environment.} All containers run on one VM with
essentially zero network latency. Detection latencies reported here are therefore
lower than what would be observed in a genuinely distributed deployment where
log shipping, network jitter, and rule evaluation would each add measurable delay.
